\documentclass[11pt,a4paper]{article}

\usepackage[margin=2.5cm]{geometry}
\usepackage{fontspec}
\setmainfont{Times New Roman}
\setmonofont[Scale=0.88]{Courier New}
\usepackage{polyglossia}
\setdefaultlanguage{vietnamese}
\usepackage{microtype}
\usepackage{mathtools,amssymb,amsmath,amsthm}
\usepackage{booktabs,tabularx,array}
\usepackage{float,needspace}
\usepackage{enumitem}
\usepackage{xspace}
\usepackage[most]{tcolorbox}
\usepackage{csquotes}
\DeclareQuoteAlias{american}{vietnamese}
\usepackage[
  backend=biber,style=authoryear-comp,language=english,
  maxcitenames=2,maxbibnames=99,uniquelist=false,
  uniquename=init,giveninits=true,doi=true,url=true,isbn=false
]{biblatex}
\DeclareLanguageMapping{vietnamese}{english}
\addbibresource{equiceval_improvement.bib}
\usepackage[unicode,colorlinks=false,pdfborder={0 0 0}]{hyperref}

\theoremstyle{plain}
\newtheorem{theorem}{Định lý}[section]
\newtheorem{proposition}[theorem]{Mệnh đề}

\theoremstyle{definition}
\newtheorem{definition}[theorem]{Định nghĩa}
\newtheorem{example}[theorem]{Ví dụ}

\theoremstyle{remark}
\newtheorem{remark}[theorem]{Nhận xét}

\newtcolorbox{keybox}[1][]{
  enhanced,breakable,colback=white,colframe=black,boxrule=0.5pt,
  arc=1pt,left=4pt,right=4pt,top=3pt,bottom=3pt,
  fonttitle=\bfseries\small,#1}

\newtcolorbox{notebox}[1][]{
  enhanced,breakable,colback=white,colframe=black,boxrule=0.9pt,
  arc=0pt,left=4pt,right=4pt,top=3pt,bottom=3pt,
  fonttitle=\bfseries\small,title=Lưu ý,#1}

\setlength{\parindent}{1.5em}
\setlength{\parskip}{0.35em}
\setlength{\emergencystretch}{3em}
\clubpenalty=10000
\widowpenalty=10000
\raggedbottom
\AtBeginBibliography{\raggedright}
\renewcommand{\arraystretch}{1.18}
\setlist[itemize]{leftmargin=1.6em,itemsep=0.3em,topsep=0.3em}
\setlist[enumerate]{leftmargin=1.8em,itemsep=0.3em,topsep=0.3em}
\setcounter{tocdepth}{2}

\newcolumntype{Y}{>{\raggedright\arraybackslash}X}
\newcolumntype{L}[1]{>{\raggedright\arraybackslash}p{#1}}
\newcolumntype{C}[1]{>{\centering\arraybackslash}p{#1}}

\newcommand{\method}{\textsf{EquiCEval}\xspace}
\newcommand{\GT}{M^{\star}}
\newcommand{\LLM}{\widehat{M}}
\newcommand{\Pgt}{P^{\star}}
\newcommand{\Pllm}{\widehat{P}}

\title{\vspace{-0.8cm}
  \textbf{\method: Đánh giá thành phần mô hình tối ưu do LLM sinh ra}\\[0.4em]
  \normalsize\textit{Đề cương nghiên cứu và kế hoạch thực nghiệm}
}
\author{Kiều Văn Tuyên}
\date{Cập nhật tháng 9 năm 2026}

\begin{document}
\maketitle

\begin{abstract}
Khi một LLM dịch yêu cầu bài toán thành mô hình tối ưu tuyến tính (LP/MILP),
cách đánh giá hiện tại có hai điểm yếu: phạt oan những cách viết tương đương về mặt toán học,
và bỏ sót lỗi trong vùng miền nghiệm chưa được kiểm tra.
\method giải quyết cả hai bằng ba bước: (1)~chuẩn hóa và ánh xạ biến để nhận ra
các cách viết tương đương, (2)~truy vấn bộ giải để tìm phương án cụ thể chứng minh
mô hình sai, và (3)~xuất phiếu chẩn đoán theo từng thành phần thay vì một điểm số duy nhất.
\end{abstract}

\tableofcontents
\newpage

%%=============================================================================
\section{Ý tưởng qua một ví dụ}
%%=============================================================================

Giả sử yêu cầu thật là $0 \leq x \leq 8$, và LLM có thể viết ra ba mô hình:
\begin{itemize}
  \item $0 \leq x \leq 8$: đúng cả cách viết lẫn ý nghĩa.
  \item $0 \leq x$ và $2x \leq 16$: cách viết khác nhưng miền nghiệm giống hệt.
  \item $0 \leq x \leq 10$: cho phép $x = 9$, trái với giới hạn thật.
\end{itemize}

So chuỗi ký tự sẽ phạt oan mô hình thứ hai. Với mô hình thứ ba,
phương án $x = 9$ là \emph{phản ví dụ}: nó được LLM cho phép nhưng yêu cầu thật cấm.
\method tìm và giải thích những bằng chứng như vậy.

\begin{keybox}[title={Thuật ngữ dùng trong tài liệu}]
\begin{description}[leftmargin=1em,labelsep=0.5em,style=nextline]
  \item[Mô hình tối ưu (formulation).] Tập biến quyết định, ràng buộc và hàm mục tiêu.
  \item[Miền nghiệm khả thi.] Tập tất cả phương án thỏa mọi ràng buộc.
  \item[Bộ giải (solver).] Phần mềm tìm nghiệm tối ưu.
  \item[Phản ví dụ.] Điểm cụ thể chứng minh hai mô hình cho kết quả khác nhau.
\end{description}
\textbf{Mô hình tham chiếu} ($\GT$): mô hình chuẩn.
\textbf{Mô hình cần đánh giá} ($\LLM$): mô hình do LLM sinh ra.
\end{keybox}

%%=============================================================================
\section{Phạm vi}
%%=============================================================================

\method tập trung vào \emph{đánh giá} mô hình LP/MILP --- không xây hệ thống sinh mô hình
hay sửa lỗi tự động.

\begin{table}[H]
\centering
\caption{Ranh giới của \method.}
\label{tab:scope}
\small
\begin{tabularx}{\textwidth}{L{3.2cm} Y Y}
\toprule
 & \textbf{Trong phạm vi} & \textbf{Ngoài phạm vi} \\
\midrule
Mục tiêu & Đánh giá mô hình LP/MILP có mô hình tham chiếu & Tự động sinh mô hình \\
Đóng góp & Đối chiếu thành phần, đo sai khác, phản ví dụ & Sửa lỗi tự động \\
Dùng LLM & Sinh mô hình để đánh giá & Huấn luyện tinh chỉnh LLM \\
Kết quả & Phát hiện và định vị lỗi & Chứng minh mô hình đúng tuyệt đối \\
\bottomrule
\end{tabularx}
\end{table}

%%=============================================================================
\section{Ba vấn đề của phương pháp cũ}
%%=============================================================================

Refai và Ahmed~\parencite{refai2025component} đã đúng khi tách biến, ràng buộc và hàm mục tiêu
để xác định lỗi ở đâu, thay vì chỉ so giá trị tối ưu. Nhưng cách đo còn ba điểm yếu.

\begin{enumerate}
  \item \textbf{Phạt oan ràng buộc tương đương.}
    Hai ràng buộc $6x + 4y \leq 12$ và $3x + 2y \leq 6$ định nghĩa cùng miền nghiệm.
    So từng hệ số sẽ cho sai số khác 0 dù không có lỗi thực sự.
    Điều tương tự xảy ra khi đổi tên biến, nhân hệ số với hằng dương,
    hoặc chuyển cận biến thành ràng buộc tường minh.

  \item \textbf{Bỏ sót lỗi trong vùng chưa lấy mẫu.}
    RMSE tính trên 100 điểm ngẫu nhiên có thể gần 0 dù mô hình thiếu một ràng buộc
    quan trọng, nếu vùng sai khác nằm ở góc miền nghiệm chưa được kiểm tra.

  \item \textbf{Thống kê chưa đủ độ bền.}
    Tương quan cao giữa Cons-RMSE và độ chênh giá trị tối ưu có thể bị chi phối
    bởi vài điểm ngoại vi. Ngoài ra, công thức gốc chia cho $|z^\star|$ gây lỗi
    khi $z^\star = 0$ (ví dụ: bài Aircraft Landing).
\end{enumerate}

%%=============================================================================
\section{Nghiên cứu liên quan}
%%=============================================================================

\method lấy Refai--Ahmed làm đối tượng cải tiến trực tiếp.
Bốn công trình dưới đây cung cấp ý tưởng thiết kế và làm phương pháp đối chứng.

\begin{table}[H]
\centering
\caption{Cách \method dùng từng nghiên cứu liên quan.}
\label{tab:sota}
\footnotesize
\begin{tabularx}{\textwidth}{Y Y Y}
\toprule
\textbf{Nghiên cứu} & \textbf{Ý tưởng dùng lại} & \textbf{Vai trò trong \method} \\
\midrule
CIR \cite{lyu2026cir}
  & Biểu diễn trung gian chuẩn hóa
  & Cơ sở chuẩn hóa, giảm phạt oan \\
EquivaMap \cite{zhai2025equivamap}
  & Kiểm tra tương đương hai mô hình
  & Phương pháp đối chứng trực tiếp \\
ReLoop \cite{lian2026reloop}
  & Kiểm tra hành vi bằng nhiễu loạn tham số
  & Đối chứng với truy vấn miền nghiệm \\
Falsification \cite{li2026falsification}
  & Phép thử bác bỏ có bằng chứng
  & Đối chiếu khả năng phát hiện lỗi \\
\bottomrule
\end{tabularx}
\end{table}

EquivaMap phù hợp khi cần kiểm tra tương đương toàn diện.
Falsification phù hợp khi không có mô hình tham chiếu.
\method cần mô hình tham chiếu và tập trung vào chẩn đoán theo từng thành phần
kèm bằng chứng từ bộ giải.

%%=============================================================================
\section{Phương pháp: \method}
%%=============================================================================

\subsection{Khái niệm cơ bản}

\begin{definition}[Mô hình tối ưu]
Một mô hình gồm $(x, D, f, \mathcal{C}, \sigma)$: biến quyết định $x$, miền kiểu $D$
(liên tục hay nguyên), hàm mục tiêu $f$, tập ràng buộc $\mathcal{C}$,
và chiều tối ưu $\sigma$.
Miền nghiệm khả thi:
\[
  F(M) = \{x \in D :\ x \text{ thỏa mọi ràng buộc trong } \mathcal{C}\}.
\]
\end{definition}

Hai mô hình được so sánh trong cùng không gian biến thông qua ánh xạ $\pi^\star$
và $\widehat\pi$ vào không gian chung $\mathcal{Y}$. Đặt:
\[
  \Pgt = \pi^\star(F(\GT)) \cap B, \qquad \Pllm = \widehat\pi(F(\LLM)) \cap B,
\]
trong đó $B$ là hộp hữu hạn xác định trước để kiểm tra.

\subsection{Kiến trúc quy trình}

Sáu mô-đun xử lý tuần tự:

\begin{equation*}
  \underbrace{\text{M0}}_{\text{Kiểm tra}}
  \to\underbrace{\text{M1}}_{\text{Chuẩn hóa}}
  \to\underbrace{\text{M2}}_{\text{Ghép cặp}}
  \to\underbrace{\text{M3}}_{\text{Đo hành vi}}
  \to\underbrace{\text{M4}}_{\text{Phản ví dụ}}
  \to\underbrace{\text{M5}}_{\text{Mục tiêu}}
\end{equation*}

M0 và M1 kiểm tra điều kiện đầu vào.
M4 tìm phản ví dụ bằng bộ giải, độc lập với kết quả ghép cặp ở M2--M3.

\subsection{M0: Kiểm tra đầu vào}

Trước khi xử lý, M0 xác định mô hình có đọc được, ánh xạ biến có thành công,
và có nghiệm hay không. Kết quả lưu trong bộ trạng thái:
\[
  S = (S_{\mathrm{parse}},\; S_{\mathrm{map}},\; S_{\mathrm{feas}},\;
       S_{\mathrm{bound}},\; S_{\mathrm{solve}}).
\]

\begin{table}[H]
\centering
\caption{Các trạng thái của M0.}
\label{tab:states}
\small
\begin{tabularx}{\textwidth}{L{3.2cm} L{3.5cm} Y}
\toprule
\textbf{Trạng thái} & \textbf{Giá trị} & \textbf{Hành động} \\
\midrule
$S_{\mathrm{parse}}$ & Đọc được / Lỗi cú pháp
  & Lỗi: dừng, báo ngay \\
$S_{\mathrm{map}}$ & Thành công / Thất bại / Không hỗ trợ
  & Thất bại: bỏ qua M4, báo chỉ số có điều kiện \\
$S_{\mathrm{feas}}$ & Có nghiệm / Vô nghiệm / Chưa biết
  & Vô nghiệm trong khi tham chiếu có nghiệm: báo lỗi ngay \\
$S_{\mathrm{bound}}$ & Bị chặn / Không giới hạn
  & Không giới hạn: chỉ truy vấn trong vùng $B$ \\
$S_{\mathrm{solve}}$ & Đã xác nhận / Hết thời gian / Lỗi số
  & Hết thời gian: báo ``chưa kết luận'' \\
\bottomrule
\end{tabularx}
\end{table}

\subsection{M1: Chuẩn hóa}

M1 đưa các cách viết tương đương về cùng một dạng để tránh phạt oan.
Mỗi ràng buộc được xử lý qua các bước:
\begin{enumerate}
  \item Đối chiếu ý nghĩa và đơn vị biến.
  \item Chuyển về dạng $g(x) \leq 0$ hoặc $h(x) = 0$.
  \item Chuyển hằng số về cùng một phía.
  \item Chia cho $q_c = \max\{\|a_c\|_1, |b_c|\}$ để chuẩn hóa tỉ lệ.
  \item Sắp xếp biến theo thứ tự từ điển.
\end{enumerate}

\begin{example}[Chuẩn hóa loại phạt oan]
$c_1: 6x + 4y \leq 12$ và $c_2: 3x + 2y \leq 6$.
Sau chuẩn hóa: $c_1 \Rightarrow \tfrac{1}{2}x + \tfrac{1}{3}y \leq 1$
và $c_2 \Rightarrow \tfrac{1}{2}x + \tfrac{1}{3}y \leq 1$.
Hai ràng buộc giống hệt nhau --- không phạt oan.
\end{example}

\subsection{M2: Ghép cặp ràng buộc}

M2 xác định ràng buộc nào của mô hình tham chiếu tương ứng với ràng buộc nào
của mô hình cần đánh giá, và gán nhãn cho từng cặp.

\begin{table}[H]
\centering
\caption{Nhãn khi đối chiếu ràng buộc.}
\label{tab:labels}
\small
\begin{tabularx}{\textwidth}{L{4.5cm} Y}
\toprule
\textbf{Nhãn} & \textbf{Ý nghĩa} \\
\midrule
Giống hệt & Biểu diễn giống nhau sau chuẩn hóa. \\
Tương đương trong ngữ cảnh & Bộ giải xác nhận tương đương khi có các ràng buộc còn lại. \\
Tương đương theo nhóm & Một nhóm ràng buộc bên này tương đương với nhóm bên kia. \\
Có phản ví dụ & Bộ giải tìm được điểm chứng minh khác nhau. \\
Chưa kết luận & Bộ giải hết thời gian hoặc chưa đủ bằng chứng. \\
\bottomrule
\end{tabularx}
\end{table}

Quy trình ưu tiên kiểm tra đơn giản trước, chỉ gọi bộ giải khi cần.
Ràng buộc dư (suy ra từ ràng buộc khác) không bị coi là lỗi.

Để kiểm tra hai ràng buộc $c_i^\star$ và $\hat{c}_j$ có tương đương không,
\method giải bài toán tìm điểm vi phạm tối đa theo mỗi chiều:
\begin{equation}
  \label{eq:eta}
  \eta_{i \to j} = \max_{\substack{y \in B \\ y\;\text{thỏa}\;c_i^\star}}
  v_{\hat{c}_j}(y), \qquad
  \eta_{j \to i} = \max_{\substack{y \in B \\ y\;\text{thỏa}\;\hat{c}_j}}
  v_{c_i^\star}(y),
\end{equation}
trong đó $v_c(y)$ là mức vi phạm chuẩn hóa của ràng buộc $c$ tại điểm $y$:
\begin{equation}
  v_c(y) =
  \begin{cases}
    \max(0, g_c(y)) / q_c, & \text{nếu ràng buộc là } g_c(y) \leq 0, \\
    |g_c(y)| / q_c,        & \text{nếu ràng buộc là } g_c(y) = 0.
  \end{cases}
\end{equation}

Ví dụ: $x \leq 8$ với $x = 10$ cho $v = (10 - 8)/8 = 0.25$.
Nếu cả $\eta_{i \to j}$ và $\eta_{j \to i}$ đều nhỏ hơn ngưỡng $\tau$,
ghi \emph{tương đương trong ngữ cảnh}; ngược lại ghi phản ví dụ.

\subsection{M3: Đo sai lệch hành vi}

M3 kiểm tra nhanh bằng $N$ điểm mẫu trong $B$, tính hai chỉ số mô tả:
\begin{equation}
  D_{\mathrm{sat}} =
  \frac{1}{N}\sum_{i=1}^{N}
  \mathbf{1}\!\left[\operatorname{sat}_{c^\star}(y^{(i)})
  \neq \operatorname{sat}_{\hat{c}}(y^{(i)})\right],
\end{equation}
\begin{equation}
  D_{\mathrm{margin}} =
  \sqrt{\frac{1}{N}\sum_{i=1}^{N}
  \bigl(v_{c^\star}(y^{(i)}) - v_{\hat{c}}(y^{(i)})\bigr)^2}.
\end{equation}

$D_{\mathrm{sat}}$ là tỷ lệ điểm mà hai ràng buộc cho kết quả trái chiều;
$D_{\mathrm{margin}}$ là mức vi phạm chênh lệch trung bình.
Đây là chỉ số mô tả bổ sung, không thay thế kết luận từ bộ giải.

\subsection{M4: Tìm phản ví dụ bằng bộ giải}

Thay vì kiểm tra điểm đã chọn sẵn, M4 hỏi bộ giải:
``Có phương án nào được bên này cho phép nhưng bên kia cấm không?''
Bộ giải kiểm tra cả hai chiều:

\begin{definition}[Hai hướng sai khác]
\label{def:delta}
\textit{Hướng $\to$} --- tìm phương án mà $\LLM$ cho phép nhưng $\GT$ cấm:
\begin{equation}
  \Delta_{\to} = \max_{c_i^\star \in \mathcal{C}^\star}
  \max_{y \in \Pllm}\, v_{c_i^\star}(y).
\end{equation}

\textit{Hướng $\leftarrow$} --- tìm phương án mà $\GT$ cho phép nhưng $\LLM$ cấm:
\begin{equation}
  \Delta_{\leftarrow} = \max_{\hat{c}_j \in \widehat{\mathcal{C}}}
  \max_{y \in \Pgt}\, v_{\hat{c}_j}(y).
\end{equation}
\end{definition}

\begin{table}[H]
\centering
\caption{Đọc kết quả hai hướng.}
\label{tab:diagnosis}
\small
\begin{tabularx}{\textwidth}{C{2.8cm} C{2.8cm} Y}
\toprule
$\Delta_{\to}$ & $\Delta_{\leftarrow}$ & Kết luận \\
\midrule
$\leq \tau$ & $\leq \tau$ & Không tìm được vi phạm đáng kể ở cả hai chiều. \\
$> \tau$ & $\leq \tau$ & $\LLM$ cho phép phương án mà $\GT$ cấm (miền rộng hơn). \\
$\leq \tau$ & $> \tau$ & $\GT$ cho phép phương án mà $\LLM$ cấm (miền hẹp hơn). \\
$> \tau$ & $> \tau$ & Sai khác cả hai chiều. \\
\bottomrule
\end{tabularx}
\end{table}

\begin{example}[LLM bỏ sót $x \leq 8$]
$\GT$: $x + y \leq 10,\; x \leq 8,\; y \leq 8$.
$\LLM$: $x + y \leq 10,\; y \leq 8$.

Bộ giải tìm điểm $(10, 0)$ trong $\Pllm$ vi phạm $x \leq 8$:
$v = (10 - 8)/8 = 0.25 > \tau \Rightarrow \Delta_{\to} = 0.25$.
\end{example}

\begin{example}[LLM siết cận thành $x \leq 5$]
$\GT$: $x + y \leq 10,\; x \leq 8$.
$\LLM$: $x + y \leq 10,\; x \leq 5$.

Bộ giải tìm điểm $(8, 0)$ trong $\Pgt$ vi phạm $x \leq 5$:
$v = (8 - 5)/5 = 0.6 > \tau \Rightarrow \Delta_{\leftarrow} = 0.6$.
\end{example}

\subsection{M5: Đánh giá hàm mục tiêu}

Hai mô hình có thể cùng miền nghiệm nhưng chọn phương án tốt nhất khác nhau.
M5 kiểm tra ba câu hỏi:
\begin{enumerate}
  \item Hệ số có đúng sau quy đổi đơn vị?
  \item Giá trị hàm mục tiêu chênh lệch tối đa bao nhiêu trên miền chung?
  \item Hai mô hình có chọn cùng phương án tối ưu không?
\end{enumerate}

Tiêu chí chính của bài là \emph{tương đương affine dương}: sau khi đưa hai bài
toán về cùng chiều tối ưu, tồn tại $a>0$ và $b$ sao cho
\begin{equation}
  f^\star(y) = a\,\tilde{f}(y) + b \qquad \text{với mọi } y.
\end{equation}
Tiêu chí này chấp nhận đổi đơn vị, nhân toàn bộ mục tiêu với một hệ số dương và
cộng một mốc cố định, vì chúng không thay đổi thứ tự tốt--xấu. Chúng tôi báo thêm
hai phân tích phụ: (i) giá trị mục tiêu bằng tuyệt đối; và (ii) chỉ cùng tập
nghiệm tối ưu. Phân tích (ii) rộng hơn: hai hàm không affine vẫn có thể tình cờ
chọn cùng một nghiệm trên một instance.

Đặt $P_\cap = \Pgt \cap \Pllm$, $f^\star$ và $\tilde{f}$ là hai hàm sau đổi đơn vị.

\textbf{Sai lệch giá trị:}
\begin{equation}
  \Delta_{\mathrm{value}} = \max_{y \in P_\cap} \frac{|f^\star(y) - \tilde{f}(y)|}{s_f},
  \qquad s_f = \max\{1,\; \|a_{f^\star}\|_1,\; |b_{f^\star}|\}.
\end{equation}

\textbf{Thiệt hại quyết định chéo} --- đo mức độ thiệt khi dùng nghiệm của mô hình
này thay cho mô hình kia:
\begin{equation}
  R_{\star \leftarrow \LLM} = \max_{y \in \widehat{A}}
  \frac{f^\star(y) - z^\star_\cap}{s_f}, \qquad
  R_{\LLM \leftarrow \star} = \max_{y \in A^\star}
  \frac{\tilde{f}(y) - \tilde{z}_\cap}{s_f},
\end{equation}
với $A^\star = \arg\min_{y \in P_\cap} f^\star(y)$,
$\widehat{A} = \arg\min_{y \in P_\cap} \tilde{f}(y)$.
Cả hai bằng 0 khi và chỉ khi hai mô hình cùng tập phương án tối ưu.

\begin{example}[Hai hàm chọn nghiệm trái ngược]
$f^\star = 3x + 2y$, $\tilde{f} = 2x + 3y$, miền $P_\cap: x + y = 6,\; x, y \geq 0$.

$f^\star$ tốt nhất tại $(0, 6)$; $\tilde{f}$ tốt nhất tại $(6, 0)$.
$R_{\star \leftarrow \LLM} = (18 - 12)/5 = 1.2$ --- hai hàm chọn nghiệm hoàn toàn trái ngược.
\end{example}

\subsection{Đầu ra: phiếu chẩn đoán}

Thay vì một điểm số duy nhất, \method trả phiếu chẩn đoán gồm:
\begin{equation}
  \begin{aligned}
  \mathcal{E}(\GT, \LLM) = \bigl(
    &S,\;\mathrm{MapCoverage},\;\mathrm{MatchCoverage},\\
    &D_{\mathrm{sat}},\; D_{\mathrm{margin}},\;
    \Delta_{\to},\; \Delta_{\leftarrow},\; \Delta_{\mathrm{value}},\\
    &R_{\star \leftarrow \LLM},\; R_{\LLM \leftarrow \star},\;
    \mathrm{Gap}_{\mathrm{sym}},\; Q_{\mathrm{total}}
  \bigr).
  \end{aligned}
\end{equation}

$\mathrm{MapCoverage}$: tỷ lệ biến tham chiếu có ánh xạ thành công.
$\mathrm{MatchCoverage}$: tỷ lệ ràng buộc tham chiếu được đối chiếu có chứng nhận.
$Q_{\mathrm{total}}$: tổng số lần gọi bộ giải.

%%=============================================================================
\section{Câu hỏi nghiên cứu}
%%=============================================================================

\begin{description}[leftmargin=1.8em,labelsep=0.5em,style=nextline]
  \item[\textbf{RQ1 --- Giảm phạt oan.}]
  \method có giảm tỷ lệ phạt oan trên các mô hình tương đương so với phương pháp gốc?

  \item[\textbf{RQ2 --- Phát hiện lỗi.}]
  Truy vấn có hướng có phát hiện lỗi miền nghiệm tốt hơn lấy mẫu ngẫu nhiên,
  đặc biệt khi giá trị tối ưu không đổi?

  \item[\textbf{RQ3 --- Ý nghĩa từng chỉ số.}]
  Từng chỉ số liên hệ thế nào với nhãn lỗi và chất lượng quyết định?
  Mối liên hệ đó có ổn định khi đổi họ bài toán?

  \item[\textbf{RQ4 --- Thứ hạng LLM.}]
  Cách đánh giá có thay đổi kết luận về LLM nào sinh mô hình chính xác hơn?

  \item[\textbf{RQ5 --- Chi phí và hiệu quả.}]
  Truy vấn có hướng tốn thêm bao nhiêu, và loại lỗi nào thì đáng bỏ chi phí đó?
\end{description}

RQ1, RQ2 và RQ5 là chính; RQ3 và RQ4 bổ sung khi đủ dữ liệu.

%%=============================================================================
\section{Bộ dữ liệu đánh giá}
%%=============================================================================

\subsection{Cấu trúc ba tập}

\begin{table}[H]
\centering
\caption{Cấu trúc bộ dữ liệu đánh giá dự kiến.}
\label{tab:benchmark}
\small
\begin{tabularx}{\textwidth}{L{3.2cm} C{1.8cm} C{2cm} Y L{2.5cm}}
\toprule
\textbf{Tập} & \textbf{Thử nhỏ} & \textbf{Đầy đủ} & \textbf{Mục tiêu} & \textbf{Nguồn} \\
\midrule
Cặp tương đương & 100--150 & 400--600 & RQ1 & Biến đổi được chứng nhận \\
Lỗi có kiểm soát & 250--400 & 800--1200 & RQ2, RQ5 & Phép sửa có nhãn \\
Đầu ra LLM thực & 100--150 & 400--600 & RQ3, RQ4 & Nhiều LLM và cách đặt câu lệnh \\
\bottomrule
\end{tabularx}
\end{table}

Mục tiêu: 50--100 bài toán gốc thuộc ít nhất 10 họ bài toán.

\subsection{Hai nhóm trọng tâm}

\textbf{Nhóm A --- Khác cách viết, cùng ý nghĩa.}
Ví dụ $x \leq 8$ và $2x \leq 16$, hoặc thêm ràng buộc được suy ra.
Công cụ chấm không nên phạt những cặp này.

\textbf{Nhóm B --- Cùng giá trị tối ưu, khác miền nghiệm.}
Ví dụ cùng tối thiểu hóa $x$: $\GT$ có $0 \leq x \leq 8$, $\LLM$ có $0 \leq x \leq 10$.
Cả hai đạt tối ưu 0 tại $x = 0$, nhưng $x = 9$ chỉ được $\LLM$ cho phép.
So giá trị tối ưu không phân biệt được; truy vấn miền nghiệm phát hiện được.

\subsection{Xây dựng dữ liệu}

\textbf{Biến đổi tạo cặp tương đương:}
\begin{itemize}
  \item Đổi tên biến và hoán vị chỉ số.
  \item Nhân ràng buộc với hằng số dương.
  \item Chuyển cận biến thành ràng buộc tường minh.
  \item Gộp/tách ràng buộc khi đã chứng minh tương đương.
  \item Thêm ràng buộc dư được suy ra.
\end{itemize}

\textbf{Phép sửa tạo lỗi kiểm thử:}
\begin{itemize}
  \item Bỏ sót ràng buộc hoặc hệ số.
  \item Đổi chiều bất đẳng thức hoặc đổi dấu.
  \item Sai hệ số, hằng số hoặc đơn vị.
  \item Sai kiểu biến (nhị phân thay liên tục).
  \item Sai tập chỉ số hoặc công thức Big-$M$.
\end{itemize}

Nhãn tham chiếu được xác minh bởi hai người đánh giá độc lập.
Báo tỷ lệ đồng thuận và hệ số Cohen $\kappa$.

%%=============================================================================
\section{Thiết kế thực nghiệm}
%%=============================================================================

\subsection{Dùng EquivaFormulation làm dữ liệu thực nghiệm}

EquivaFormulation là bộ dữ liệu gồm các bài toán tối ưu gốc kèm nhiều biến thể
được tạo ra theo các quy tắc đã định. Chúng tôi dùng đây làm nguồn dữ liệu thực nghiệm,
còn EquivaMap \parencite{zhai2025equivamap} được dùng làm phương pháp đối chứng.

\textbf{Ba quy tắc để dùng dữ liệu này một cách công bằng:}

\textbf{(D1) Ánh xạ tên biến được cấp tường minh, không phải đáp án.}
EquivaFormulation cung cấp kèm file \texttt{variable\_mappings.json} xác định tương ứng
giữa tên biến dài trong bài gốc (ví dụ: \texttt{NumberOfCars}) và tên biến ngắn
trong biến thể (ví dụ: \texttt{r}).
File này chỉ đề xuất hai hệ tọa độ cần so sánh; nó không chứng minh hai mô hình
tương đương. Trong thí nghiệm ``có mapping'', file được cấp làm đầu vào. Trong
thí nghiệm ``không mapping'', công cụ chỉ dùng phép dò một--một bảo thủ. So sánh
này đo ích lợi của mapping; không được dùng để khẳng định công cụ đã hiểu ngữ nghĩa
của mọi tên biến.

\textbf{(D2) Phân biệt rõ ``Labeled'' và ``Full''.}
Một cặp chỉ thuộc tập \textbf{Labeled} khi nhãn theo tiêu chí của bài này được
xác minh độc lập. \textbf{Full} là toàn bộ cặp đọc được từ nguồn dữ liệu, gồm cả
cặp chưa có nhãn hoặc chưa hỗ trợ. Vì vậy, ``Labeled'' không có nghĩa đơn giản là
``bỏ đúng hai loại biến thể''. Ánh xạ nhiều--một, biến phụ nguyên, hoặc candidate
không có phép đổi tọa độ một--một có thể phải để \emph{chưa kết luận}. FPR và Recall
chỉ dùng Labeled; Full chỉ báo độ phủ và số cặp chưa xử lý.

\textbf{(D3) Tách ba ý nghĩa của hàm mục tiêu.}
Tiêu chí chính chấp nhận $f(x) \Leftrightarrow k f(x)+b$ với $k>0$, sau khi đưa
min/max về cùng chiều; do đó phép nhân toàn bộ mục tiêu với hằng số dương không
phải lỗi. Bảng phụ ``giá trị tuyệt đối'' vẫn phân biệt hai hàm khi con số mục tiêu
khác nhau. Bảng phụ ``argmin'' chỉ hỏi chúng có chọn cùng tập nghiệm tối ưu hay
không. Mỗi bảng dùng tệp nhãn được tạo riêng theo đúng tiêu chí; không dùng nhãn
của bảng này để chấm bảng khác.

\textbf{Kiểm toán trước khi báo số liệu.}
Các kết quả tạo trước khi áp dụng ba quy tắc trên chỉ là nhật ký để đối chiếu, không
phải bằng chứng trong bài. Pipeline hiện tại không sao chép cận của candidate sang
reference, không thay vô cực bằng một cận lớn, không khử tùy tiện biến phụ nguyên,
và phát lại phản ví dụ trên mô hình gốc. Bảng FPR/Recall chỉ được lập lại sau khi
toàn bộ cặp đã được relabel và chạy lại bằng phiên bản này.

\subsection{Các phương pháp đối chứng}

\begin{enumerate}
  \item Bản tái triển khai Refai--Ahmed~\parencite{refai2025component}.
  \item So khớp biểu diễn giống hệt (exact match).
  \item Chuẩn hóa nhưng không dùng bộ giải kiểm tra ngữ cảnh.
  \item EquivaMap~\parencite{zhai2025equivamap} trên các trường hợp tương thích.
  \item Nhiễu loạn tham số (kiểu ReLoop~\parencite{lian2026reloop}).
  \item \method đầy đủ.
\end{enumerate}

\subsection{Chỉ số đánh giá}

\textbf{Tỷ lệ phạt oan} (FPR): trong các cặp thật sự tương đương,
bao nhiêu cặp bị báo nhầm là sai?

\textbf{Tỷ lệ phát hiện} (Recall): trong các cặp thật sự sai,
bao nhiêu cặp được phát hiện?

\[
\begin{aligned}
 \mathrm{FPR}
   &= \frac{\#\text{cặp tương đương bị báo nhầm}}
           {\#\text{cặp tương đương có nhãn xác minh}}, \\[0.5em]
 \mathrm{Recall}
   &= \frac{\#\text{cặp sai được phát hiện}}
           {\#\text{cặp sai có nhãn xác minh}}.
\end{aligned}
\]

FPR thấp tốt; Recall cao tốt. Phải báo cả hai cùng nhau.
Mẫu chưa kết luận không tính là phát hiện, cũng không tính là tương đương.

Các chỉ số bổ sung:
\begin{itemize}
  \item Độ chính xác ánh xạ biến.
  \item Độ chính xác định vị lỗi (thành phần báo đầu tiên có đúng không).
  \item Tỷ lệ phản ví dụ qua kiểm tra độc lập.
  \item Thời gian chạy, số lần gọi bộ giải, tỷ lệ chưa kết luận.
\end{itemize}

\subsection{Thử bỏ từng mô-đun}

\begin{table}[H]
\centering
\caption{Kế hoạch thử bỏ từng mô-đun.}
\label{tab:ablation}
\small
\begin{tabularx}{\textwidth}{L{3.5cm} Y Y}
\toprule
\textbf{Bỏ đi} & \textbf{Câu hỏi} & \textbf{Chỉ số nhạy nhất} \\
\midrule
Chuẩn hóa & Có phạt oan lại không? & FPR trên tập tương đương \\
Ánh xạ biến & Bỏ sót hoặc ghép sai khi đổi tên? & Độ chính xác ánh xạ \\
Kiểm tra ngữ cảnh & Ràng buộc dư có bị phạt không? & FPR, độ chính xác đối chiếu \\
Truy vấn có hướng & Lấy mẫu ngẫu nhiên có bỏ sót không? & Recall, độ phủ phản ví dụ \\
\bottomrule
\end{tabularx}
\end{table}

%%=============================================================================
\section{Đóng góp và kết quả mục tiêu}
%%=============================================================================

Sau khi hoàn thành thực nghiệm, bài báo kỳ vọng đóng góp:

\begin{enumerate}
  \item Một quy trình đánh giá thành phần kèm bằng chứng rõ ràng:
    ánh xạ biến, kiểm tra từng ràng buộc, phản ví dụ có thể kiểm tra lại.
  \item Một bộ dữ liệu tập trung vào hai điểm mù: cách viết khác nhưng cùng nghĩa,
    và cùng giá trị tối ưu nhưng khác miền nghiệm.
  \item Nghiên cứu thực nghiệm chỉ ra khi nào truy vấn có hướng hữu ích hơn
    lấy mẫu ngẫu nhiên và khi nào không.
\end{enumerate}

Bài báo được coi là thuyết phục khi:
\begin{itemize}
  \item Giảm rõ FPR trên tập cặp tương đương.
  \item Tăng Recall trên lỗi bỏ sót và lỗi siết cận.
  \item Phân tích đánh đổi chi phí--khả năng phát hiện trên cùng ngân sách tính toán.
\end{itemize}

%%=============================================================================
\section{Kế hoạch 12 tuần}
\label{sec:plan}
%%=============================================================================

\begin{description}[leftmargin=2.5em,labelsep=0.5em,style=nextline]
  \item[\textbf{Tuần 1--2.}]
  Xác định rõ tiêu chí tương đương; viết bộ kiểm tra với đáp án biết trước.
  \item[\textbf{Tuần 3--4.}]
  Hoàn thiện chuẩn hóa và ánh xạ biến; tái triển khai phương pháp đối chứng;
  thiết lập nhật ký chạy cho từng cặp.
  \item[\textbf{Tuần 5--6.}]
  Hoàn thiện hai hướng truy vấn và kiểm tra phản ví dụ độc lập; đo ngân sách thực.
  \item[\textbf{Tuần 7--8.}]
  Tạo đợt thử nhỏ; xác nhận nhãn độc lập; mở rộng sau khi nhãn đã ổn định.
  \item[\textbf{Tuần 9.}]
  Thu khoảng 120 đầu ra LLM thực (2 LLM, 2 cách đặt câu lệnh, 3 lần sinh, 10 bài toán).
  \item[\textbf{Tuần 10.}]
  Chạy phương pháp đối chứng, thử bỏ từng mô-đun.
  \item[\textbf{Tuần 11.}]
  Tính khoảng tin cậy (bootstrap theo họ bài toán), rà soát phản ví dụ.
  \item[\textbf{Tuần 12.}]
  Viết kết quả từ nhật ký chạy; công bố cấu hình và hướng dẫn tái lập.
\end{description}

\subsection{Rủi ro và cách xử lý}

\begin{table}[H]
\centering
\caption{Rủi ro chính và hành động.}
\label{tab:risk}
\small
\begin{tabularx}{\textwidth}{L{3cm} L{2.8cm} Y}
\toprule
\textbf{Rủi ro} & \textbf{Dấu hiệu sớm} & \textbf{Hành động} \\
\midrule
Ánh xạ thất bại nhiều
  & $>$20\% ca thử thất bại
  & Chỉ nhận ánh xạ tường minh; báo không hỗ trợ \\
MILP quá chậm
  & $>$30\% hết thời gian
  & LP làm nghiên cứu chính; MILP là mở rộng \\
Lỗi đọc đầu ra LLM
  & $>$10\% không đọc được
  & Báo tỷ lệ toàn quy trình riêng \\
Không cải thiện được baseline
  & Recall không tăng trên nhóm lỗi chính
  & Báo kết quả và phân tích giới hạn \\
\bottomrule
\end{tabularx}
\end{table}

%%=============================================================================
\section{Cấu trúc bài báo dự kiến}
%%=============================================================================

\begin{enumerate}
  \item \textbf{Giới thiệu:} tại sao so giá trị tối ưu không đủ;
    tại sao đánh giá theo thành phần cần cải tiến.
  \item \textbf{Nghiên cứu liên quan:} EquivaMap, Falsification, CIR, ReLoop.
  \item \textbf{Phân tích điểm yếu:} ba điểm mù với ví dụ cụ thể.
  \item \textbf{\method:} ánh xạ biến, chuẩn hóa, đối chiếu, truy vấn, mục tiêu.
  \item \textbf{Bộ dữ liệu:} cách xây dựng và xác minh nhãn.
  \item \textbf{Thực nghiệm:} RQ1, RQ2, RQ5 là chính; RQ3--RQ4 bổ sung.
  \item \textbf{Phân tích:} phản ví dụ, chi phí theo Recall, thử bỏ từng mô-đun.
  \item \textbf{Giới hạn:} cần mô hình tham chiếu; ánh xạ tùy ý phải cấp tường minh.
  \item \textbf{Kết luận:} từ đếm thành phần sang chẩn đoán có phản ví dụ.
\end{enumerate}

%%=============================================================================
\Needspace{12\baselineskip}
\section{Thứ tự ưu tiên}
%%=============================================================================

\begin{keybox}[unbreakable,title=\textbf{Làm đúng thứ tự này}]
\begin{enumerate}
  \item Kiểm tra đúng trên những ca nhỏ có đáp án biết trước.
  \item Tạo nhãn đã xác minh độc lập, đặc biệt cho lỗi giữ nguyên giá trị tối ưu.
  \item Kiểm tra trên nhiều họ bài toán mới, không chỉ tăng biến thể của vài bài cũ.
  \item Đo bằng nhật ký chạy thật; báo cả chi phí và những ca chưa kết luận.
  \item Viết khẳng định mạnh đến mức bằng chứng thực tế cho phép.
\end{enumerate}
\end{keybox}

\printbibliography[title={Tài liệu tham khảo}]
\end{document}
